\section{Scheduling Other Types of DL Workloads}
\label{sec_other_workload}




The previous two sections categorize the scheduling of general training and inference workloads, respectively. In this section, we discuss the works for optimizing other types of DL workloads.
Table~\ref{tab_other} presents the past works for hyperparameter optimization workloads as well as hybrid training and inference workloads, with the detailed description as below. 

\subsection{Hyperparameter Optimization Workloads}

A Hyperparameter Optimization (HPO) job aims to identify the best hyperparameters for a DL task. Technically speaking, HPO belongs to the category of DL training workloads. Here we discuss it separately because it has unique features compared to the general DL training jobs. Specifically, a HPO job typically needs to search a set of hyperparameter configurations. Each configuration is associated with a \emph{trial} \cite{RubberBand}, which is a common DL training workload containing training and evaluation procedures. However, in the HPO context, these trails are extremely similar and orchestrated by a HPO scheduling policy. To accelerate the HPO process, the policy can kill poor-performing trails through early stopping and allocate more resources to promising trails. These optimizations on HPO workloads can deliver remarkable acceleration.





\textbf{Accuracy efficiency.}
Hermes~\cite{Hermes} is a scheduler to expedite HPO workloads in GPU datacenters. It provides a container preemption mechanism to enable migration between DL jobs with minimal overhead. Besides, it also considers the algorithmic property of HPO workloads and devises a convergence-aware scheduling algorithm to favor promising hyperparameter configurations.
HyperSched~\cite{HyperSched} aims to boost the accuracy performance of HPO workloads within the given time and resource budgets. In an HPO workload, the promising hyperparameter trial generally has higher accuracy at the early training stage. Then, HyperSched allocates more resources to the promising trials and terminates unpromising ones. With this technique, HyperSched can parallel orders of magnitude more hyperparameter trials and therefore maximize the final accuracy substantially. As an extension of HyperSched~\cite{HyperSched}, SEER~\cite{SEER} replaces the resource budget with the monetary cost budget in the cloud. Without the resource constraint, SEER can launch numerous training trials at the early training stage to explore promising hyperparameter combinations. Then it will terminate poor training trials and allocate more cost budget to promising trials during the training progress. The accommodation of more trials at the hyperparameter exploration stage can improve the final model accuracy. HyperDrive~\cite{HyperDrive} is another scheduler framework for hyperparameter exploration as well.
It develops an accuracy curve fitting model to extrapolate the accuracy in the subsequent training iterations. The accuracy prediction engine allows HyperDrive to terminate low-quality jobs early and adjust the resource allocation dynamically. Moreover, Fluid \cite{Fluid} further leverages the job packing mechanism to improve resource utilization and accelerate the HPO process.


\textbf{Cost efficiency.}
RubberBand~\cite{RubberBand} aims to reduce the monetary cost of an HPO job with the constraint of timing budget in the cloud. Since the optimal amount of resources for an HPO workload differs in early and later stages of the hyperparameter search and model optimization processes, RubberBand needs to jointly accommodate the job throughput, resource amount and cloud pricing. By building a profiling model, RubberBand can predict the job throughput and corresponding cost for a given resource allocation policy. Then it uses this model to generate an optimal cost-efficient solution for satisfactory cost reduction.

\begin{table*}[t]
    \centering
    \caption{\textbf{Summary of schedulers for other workloads in GPU Datacenters.}}
    \vspace{-10pt}
    \resizebox{\textwidth}{!}{
        \begin{tabular}{@{}cccccccc@{}}
            \toprule
            \textbf{Type}           & \textbf{Year}               & \textbf{Scheduler}                        & \textbf{Objectives}                                                                                                                                                             & \textbf{Approaches}                           & \textbf{Advantages}                         & \textbf{\makecell[c]{Exp.             \\ Scale}} & \textbf{\makecell[c]{Source \\ Code}} \\ \midrule
            \multirow{6}{*}{HPO}    & 2017                        & HyperDrive \cite{HyperDrive}              & \ding{168}                                                                                                                                                                      & Dynamic Probabilistic Accuracy Prediction     & JCT Reduction                               & S                         & -         \\
                                    & 2019                        & HyperSched \cite{HyperSched}              & \ding{171}                                                                                                                                                                      & ASHA; Dynamic Resource Allocation             & Efficient HPO under DDL                     & S                         & \ding{52} \\
                                    & 2021                        & Hermes \cite{Hermes}                      & \ding{168}\ding{169}                                                                                                                                                            & Time-sharing Execution with Low Overhead      & JCT Reduction                               & S                         & -         \\
                                    & 2021                        & RubberBand \cite{RubberBand}              & \ding{168}\ding{169}                                                                                                                                                            & Model JCT and Cost Prior to Runtime           & Cost-Effectiveness                          & S                         & -         \\
                                    & 2021                        & SEER \cite{SEER}                          & \ding{168}\ding{169}                                                                                                                                                            & Dynamic Resource Allocation                   & Cost Constraint HPO                         & S                         & -         \\
                                    & 2021                        & Fluid \cite{Fluid}                        & \ding{95}                                                                                                                                                  \ding{168}\ding{169} & Job Packing; Elastic
            Training                & Better Resource Utilization & S                                         & \ding{52}                                                                                                                                                                                                                                                                                                             \\ \midrule
            \multirow{4}{*}{Hybrid} & 2018                        & Rafiki \cite{wang2018rafiki}              & \ding{170}\ding{169}                                                                                                                                                            & RL to optimize accuracy and latency           & Unified platform for training and inference & S                         & \ding{52} \\
                                    & 2019                        & Kube-Knots \cite{Thinakaran2019kubeknots} & \ding{169}\ding{95}                                                                                                                                                             & Dynamic container resizing                    & GPU utilization-aware colocation            & S                         & -         \\
                                    & 2020                        & CMS \cite{li2020cms}                      & \ding{170}\ding{95}                                                                                                                                                             & Common architecture for trainers and modelets & Continuous learning                         & S                         & -         \\
                                    & 2022                        & Aryl \cite{Aryl}                          & \ding{95}\ding{170}                                                                                                                                                             & Capacity Loaning                              & Cluster Size Extension                      & L                         & -         \\ \bottomrule
        \end{tabular}
    }

    \begin{flushleft}
        \begin{tablenotes}[para,flushleft]
            \scriptsize
            \textbf{Objectives:} \ding{95} Utilization \ding{168} Makespan \ding{169} Cost \ding{170} Accuracy \ding{171} DDL; \hspace{10pt}  \textbf{Experiment GPU Scales:} the scale of physical testbed. S $(0, 30]$ M $(30, 60]$ L $(60, 120]$ XL $(120, \infty]$ -: no evaluation on a physical cluster or not clearly specified.
        \end{tablenotes}
    \end{flushleft}

    \vspace{-10pt}
    \label{tab_other}
\end{table*}


\subsection{Hybrid of Training and Inference Workloads}
In Sec. \ref{sec_training} and Sec. \ref{sec_inference} we consider the scheduling techniques for training and inference solutions separately. Actually, there are some DL systems that host these two workloads in a unified manner. Due to the distinct features of the two types of workloads, new solutions are needed for efficient scheduling in GPU datacenters. Below we review and summarize these works.


\textbf{End-to-end DL development.}
Some works consider the entire pipeline of DL development and deployment, including model training, inference, as well as periodically retraining and updating.
CMS~\cite{li2020cms} designs a continuous machine learning and serving platform, which orchestrates the model training executions, model deployments, and model update services. It unifies different training jobs into a simple trainer contract and provides common essential pipeline abstracts of training. Then the scheduler monitors the resource consumption of these resource-intensive training tasks to avoid contention. Other techniques including model validation are also applied to guarantee the quality of newly-trained models.
Rafiki~\cite{wang2018rafiki} proposes to reuse the datasets and parameters across different training and inference jobs. It implements a unified distributed dataset storage and a parameter server. The parameter server is kept in memory and shares model parameters across different trials in the training and dumped for later reuse in the inference. Other common underlying components are also shared across training and inference workloads to reduce the operational cost, e.g., storage, communication protocols, and compute resources.







\textbf{Mix of training and inference development.}
Some works optimize the datacenter with mixed workloads of DL training and inference.
Kube-Knots~\cite{Thinakaran2019kubeknots} minimizes the resource waste by allocating offline batch jobs to better utilize the spare resources. It discovers that the GPU energy efficiency increases with its utilization. Therefore, achieving higher GPU utilization indicates higher energy efficiency and less resource wastes. To this end, Kube-Knots colocates the predictable GPU batch jobs with the online inference workloads, as the inference jobs generally underutilize the GPUs. During the scheduling, it also predicts the peak resource utilization of the jobs from their online resource usages.
Aryl~\cite{Aryl} also observes the resource can be wasted due to the resource over-provisioning for burst inference requests. It proposes the concept of capacity loaning, which allows the inference cluster to loan the idle GPU servers during low-traffic periods to run training jobs. Since the preemption cost is not negligible, it minimizes the total number of preemptions during the scheduling. Combined with the elasticity on resource demands for part of training jobs, Aryl successfully guarantees the timely resource allocation to inference jobs via a heuristic method.


\subsection{Implications}

In addition to optimizing general DL workloads, there are some opportunities for further resource efficiency improvement. One direction is to apply specific optimizations for hyperparameter search workloads. Compared with the workload-agnostic manner, it can deliver over one magnitude resource and time conservation. However, how to integrate this mechanism well into the scheduler and coordinate with general DL workloads is a challenging topic, requiring more future research investigation. Besides, considering the whole DL model development pipeline instead of focusing on a certain stage can bring extra system performance enhancement. For instance, breaking the shackles between the training and inference cluster resource not only considerably diminishes the queuing delay of training workloads but also improves the model serving quality. These directions deserve more attention in future research on GPU datacenter scheduling systems.


